\begin{summarybox}
\textbf{\textcolor{CSummary}{先看得懂的结论，再决定要不要背公式}}

\textbf{第一句话：} 在圆外一点作两条切线，\textbf{两个切点的连线}叫这个点的极线；点在圆上时，同一公式得到的是切线；点在圆内时，同一公式仍给出一条直线，但不是能实际画出来的切点弦。\textbf{这就是为什么需要引入“极线”这个统一名称。}

\medskip
\textbf{\textcolor{CSummary}{最值得记住的两个结果}}
\[
\boxed{Q\in l_P\iff P\in l_Q}
\]
意思：$Q$ 在 $P$ 对应的直线上，当且仅当 $P$ 也在 $Q$ 对应的直线上。
\[
\boxed{Q\text{ 在固定的 }l_P\text{ 上移动}\quad\Longrightarrow\quad l_Q\text{ 始终经过固定的 }P}
\]
意思：\textbf{动点在定直线上，对应的动直线经过定点}。它通常比“共轭”这个名称更有解题价值。

\medskip
\textbf{\textcolor{CSummary}{一眼能检查的圆形母题}}

圆 $x^2+y^2=4$，固定点 $P(4,0)$：
\[
\underbrace{l_P:x=1}_{\text{固定直线}},\qquad
\underbrace{Q(1,t)}_{\text{直线上的动点}},\qquad
\underbrace{l_Q:x+ty=4}_{\text{经过 }P(4,0)\text{ 的动直线}}.
\]
\textbf{核心检查：} 把 $(4,0)$ 代入 $x+ty=4$，左边恒等于 $4$。圆内、圆上、圆外都可这样验证；只有极线与圆的相交情况会随着 $t$ 变化。

\medskip
\textbf{\textcolor{CSummary}{四类曲线的对照（不必一次背完）}}
\begin{itemize}[leftmargin=*,itemsep=0.35em]
\item 圆 $x^2+y^2=r^2$：$l_P:x_1x+y_1y=r^2$。
\item 椭圆 $\dfrac{x^2}{a^2}+\dfrac{y^2}{b^2}=1$：$l_P:\dfrac{x_1x}{a^2}+\dfrac{y_1y}{b^2}=1$。
\item 双曲线 $\dfrac{x^2}{A^2}-\dfrac{y^2}{B^2}=1$：$l_P:\dfrac{x_1x}{A^2}-\dfrac{y_1y}{B^2}=1$。
\item 抛物线 $y^2=2px$：$l_P:y_1y=p(x+x_1)$。
\end{itemize}
共同点只有一个：\textbf{把两点坐标交换，检验式不变。} 理由仅仅是加法、乘法的交换律。

\medskip
\par\noindent\begin{minipage}{\linewidth}
\textbf{\textcolor{CSummary}{考场何时使用？}}
\begin{itemize}[leftmargin=*,itemsep=0.34em]
\item 看见\textbf{外点两切线、切点弦、动点在直线上、证明动线过定点}，先考虑能否认出极线。
\item \textbf{先验证固定直线确实是某个固定点的极线}，否则不能无条件套“动线过定点”。
\item 认出之后，利用 $Q\in l_P\Rightarrow P\in l_Q$ 预测定点，最后用代入法核验。
\item 若只是普通求坐标、交点或斜率，没有上述特征，优先常规解法。
\end{itemize}

\end{minipage}\par
\medskip
\textbf{\textcolor{CSummary}{三个必须排除的误解}}
\begin{itemize}[leftmargin=*,itemsep=0.3em]
\item \textbf{圆内无真实切点弦：} 不要把圆内的极线画成由圆内点引出的两条切线。
\item \textbf{中心点例外：} 对圆、椭圆、双曲线，其中心不能按这里的方法配出普通极线。
\item \textbf{不是所有固定直线都有普通极点：} 证明动线过定点之前，必须先有 $l=l_P$ 这一前提。
\end{itemize}

\medskip
\textbf{\textcolor{CSummary}{学习定位}}

这是一条\textbf{有较高理论统一价值、但不宜作为全体高中生必背高频技巧}的提高型结论。入门只要求掌握圆的例子和交换律证明；掌握切点弦过定点后，再把它迁移到椭圆、双曲线、抛物线。\textbf{术语只作为简写使用：不理解图形和方程，就没有必要先背名称。}
\end{summarybox}
